\chapter{Introdução}
\label{chap:intro}

\section{Enquadramento}
A prevalência de doenças crónicas associadas a hábitos alimentares inadequados, como a obesidade, a diabetes tipo 2 e as doenças cardiovasculares, tem aumentado a atenção dada à nutrição como fator de saúde. Neste contexto, o acompanhamento nutricional personalizado tornou-se uma ferramenta relevante na promoção de estilos de vida mais saudáveis e na melhoria dos indicadores clínicos dos pacientes.

A eficácia de uma dieta não depende apenas do plano alimentar em si. Depende também de fatores biológicos do paciente, do seu comportamento e nível de adesão, e até do perfil do profissional que faz o acompanhamento. Esta combinação de fatores torna difícil prever, à partida, que planos têm maior probabilidade de sucesso para um determinado perfil de paciente.

\section{Motivação}
O interesse deste problema está em perceber se os dados recolhidos durante o acompanhamento nutricional contêm informação suficiente para antecipar resultados. Se for possível estimar com alguma fiabilidade quanto peso um paciente vai perder, ou se uma dieta tem boa probabilidade de sucesso, então essa informação pode apoiar decisões clínicas antes de a dieta começar.

O grupo abordou o problema através de um conjunto de dados que reúne informação sobre pacientes, planos dietéticos, nutricionistas responsáveis e resultados das dietas ao fim de seis meses. O objetivo foi explorar e caracterizar esses dados, identificar padrões associados ao sucesso ou insucesso das dietas, e construir modelos capazes de fazer previsões úteis.

\section{Objetivos}
O trabalho consistiu na construção de um conjunto de \textit{scripts} em Python que executam, de forma encadeada, todas as etapas de um processo de análise de dados. As etapas vão desde a integração dos quatro ficheiros \ac{CSV} fornecidos até à aplicação de modelos de previsão e à discussão crítica dos resultados obtidos.

Em concreto, os objetivos foram integrar os dados num único conjunto, limpar e transformar esse conjunto, realizar uma análise exploratória detalhada, identificar grupos de pacientes através de \textit{clustering}, treinar modelos supervisionados para prever a perda de peso, e por fim comparar e discutir o desempenho desses modelos.

\section{Organização do Documento}
O documento está organizado em seis capítulos. Depois desta introdução, o segundo capítulo descreve a integração dos dados e a estrutura do conjunto unificado. O terceiro apresenta a limpeza e o processamento aplicados. O quarto descreve a análise exploratória e os padrões diretos encontrados. O quinto trata da identificação de padrões por \textit{clustering} e da aplicação dos modelos de previsão. O último capítulo apresenta a discussão crítica dos resultados e as conclusões do trabalho.
